% Generated draft: requires mathematical and source review.
\documentclass[11pt]{article}
\usepackage[margin=25mm]{geometry}
\usepackage{fontspec}
\setmainfont{Latin Modern Roman}
\usepackage{amsmath,amssymb,mathtools,unicode-math}
\setmathfont{Latin Modern Math}
\usepackage{enumitem,needspace,xurl,hyperref}
\setlength{\parindent}{0pt}
\setlength{\parskip}{5pt}
\setlength{\emergencystretch}{3em}
\newcommand{\N}{\mathbb{N}}
\newcommand{\Z}{\mathbb{Z}}
\newcommand{\Q}{\mathbb{Q}}
\newcommand{\R}{\mathbb{R}}
\newcommand{\C}{\mathbb{C}}
\newcommand{\F}{\mathbb{F}}
\newcommand{\A}{\mathbb{A}}
\newcommand{\PP}{\mathbb{P}}
\newcommand{\E}{\mathbb{E}}
\newcommand{\eps}{\varepsilon}
\newcommand{\dd}{\,\mathrm{d}}
\begin{document}
\section*{UnsolvedMath ID: 800549}
\textbf{Source problem code:} catalogue entry
\subsection{Definitions and mathematical statement}
The source asks for a characterization of bar $k$-visibility graphs, but does not supply the definition of this class. The following formulation therefore uses an explicitly provisional visibility convention.

Let $k\in\Z_{\geq 0}$ and let $G=(V,E)$ be a finite simple undirected graph. A bar representation assigns to each $v\in V$ a nondegenerate closed horizontal line segment $B_v\subset\R^2$, with the segments pairwise disjoint. Two distinct bars $B_u,B_v$ are called $k$-visible if there is an open axis-parallel rectangle of positive width whose lower and upper horizontal sides lie in these two bars, in either order, and whose interior intersects at most $k$ other bars. Under the provisional convention, $G$ is a bar $k$-visibility graph if it has such a representation satisfying
\[
\{u,v\}\in E
\quad\Longleftrightarrow\quad
B_u\text{ and }B_v\text{ are }k\text{-visible}
\qquad(u\ne v).
\]

\textbf{Characterization problem.} For each nonnegative integer $k$, characterize the finite graphs that are bar $k$-visibility graphs by necessary and sufficient conditions. The source does not specify the desired form of the characterization. The definition above must be checked against the cited source before this draft can be treated as a precise transcription: in particular, positive-width visibility, endpoint conventions, disjointness requirements, and the equivalence between edges and visible pairs are not specified in the supplied excerpt.

For context, with $n=|V|$, the supplied record attributes to Dean et al. the edge upper bound
\[
(k+1)\left(3n-\frac{7}{2}k-5\right)-1,
\]
and a conjectured improvement to
\[
(k+1)(3n-4k-6),
\]
which it reports was subsequently proved by Hartke et al. It also attributes the chromatic bound $\chi(G)\leq 6k+6$ to Dean et al., and mentions work on $d$-regular $k$-visibility graphs, meaning graphs in which every vertex has degree $d$. The excerpt supplies no parameter ranges for these background bounds; they are preserved as reported context, not asserted here without their original hypotheses. The record contrasts this characterization question with the stated existing characterization of bar visibility graphs.
\subsection{Short English statement}
For each nonnegative integer $k$, describe exactly which finite graphs admit a representation by horizontal bars in which adjacency corresponds to vertical visibility through at most $k$ other bars. The visibility convention stated here is provisional because the supplied excerpt omits its definition. The registry labels the problem open; this status is supplied metadata, not an independently verified assessment.
\subsection{Sources}
\begin{itemize}
\item Local UnsolvedMath registry, supplied complete source record: permanent ID $800549$, source problem code \texttt{catalogue entry}, title ``Visibility Graphs of Points, Segments and Polygons --- Problem 21.'' The record attributes the problem to Ghosh and Goswami, with proposed year 2010, and labels it open. Supplied upstream revision: \texttt{d371ee917ac9647fbb06cd6fc529deec6f3dda82}.
\item Ghosh and Goswami, \emph{Unsolved Problems in Visibility Graphs of Points, Segments and Polygons} (2010), Problem 21, \url{https://arxiv.org/abs/1012.5187}. The supplied provenance records access on 2026-10-16 and extraction by \texttt{arxiv\_latex}; these are supplied provenance claims, not evidence of independent source consultation for this draft.
\item Dean et al., source reference label \texttt{deglst-bvg-2007}, cited in the supplied excerpt for edge bounds, a chromatic bound, regular graphs, and the characterization question. Full bibliographic details are not supplied.
\item Hartke et al., source reference label \texttt{hvw-frbvg-2007}, cited in the supplied excerpt for the improved edge bound and the characterization question. Full bibliographic details are not supplied.
\end{itemize}
\end{document}
